% Self-contained report. Compile with pdfLaTeX in Overleaf (twice for contents).
\documentclass[12pt,a4paper]{article}
\usepackage[T1]{fontenc}
\usepackage[utf8]{inputenc}
\usepackage[english]{babel}
\usepackage[margin=2.5cm]{geometry}
\usepackage{xcolor}
\usepackage{listings}
\usepackage{booktabs}
\usepackage{microtype}
\usepackage[hidelinks]{hyperref}
\definecolor{listinggray}{RGB}{247,247,247}
\lstset{basicstyle=\ttfamily\footnotesize,backgroundcolor=\color{listinggray},
 frame=single,rulecolor=\color{gray},breaklines=true,breakatwhitespace=false,
 columns=fullflexible,keepspaces=true,showstringspaces=false,tabsize=4,
 captionpos=b,aboveskip=10pt,belowskip=10pt}
\setlength{\emergencystretch}{3em}
\begin{document}
\begin{titlepage}
\centering
\textbf{Ministry of Education, Culture and Research of the Republic of Moldova}\par\medskip
\textbf{Technical University of Moldova}\par\medskip
\textbf{Faculty of Computers, Informatics and Microelectronics}\par\medskip
\textbf{Department of Software Engineering and Automatic}\par
\vspace{2.5cm}
\textbf{ANDREI ELCISCEV}\par
\textbf{Group FAF-243}\par
\vspace{2cm}
{\Large\textbf{LABORATORY WORK No. 10}\par}
\medskip
{\Large\textbf{ADD YOUR OWN SYSTEM CALL: TRACE}\par}
\vspace{2cm}
\begin{flushright}
\textbf{Checked by:}\\
Rei\c{t}man Patricia
\end{flushright}
\vfill
\textbf{Chi\c{s}in\u{a}u, 2026}
\end{titlepage}
\pagenumbering{arabic}
\tableofcontents
\newpage
\section{Introduction}
This laboratory work extends the xv6 RISC-V teaching operating system with a new
system call, \texttt{trace}. A process supplies a bit mask selecting which system
calls should be printed after execution. Each selected trace line contains the
process identifier, the system call name, and its return value.

Lab 9 observed Linux through \texttt{strace}; this work implements a tracing
mechanism inside the xv6 kernel. The experiment follows syscall number 23 from
the user declaration and generated assembly stub to the kernel dispatcher and
handler, then tests mask selection and inheritance across process creation.

The report follows the structure of the previous laboratory reports. All terminal
listings are captured from the actual VM and QEMU sessions. Excerpts omit build
verbosity or unrelated commands only where identified; no screenshots or external
assets are required to compile this document.

\section{Environment and Initial Verification}
A working directory was created at \texttt{/home/user/os-lab10}. The existing
\texttt{/home/user/xv6-riscv} tree contained earlier coursework, so it was copied
with its Git history into \texttt{/home/user/os-lab10/xv6-riscv}. The original tree
was preserved. Existing coursework source changes were recorded in a baseline
commit before the Lab 10 implementation, on branch \texttt{codex/lab10-trace}.

\begin{lstlisting}[caption={Working tree, kernel, RISC-V compiler, emulator, and README size}]
$ bash -c pwd; uname -sr; riscv64-linux-gnu-gcc --version | head -1; qemu-system-riscv64 --version | head -1; wc -c README
/home/user/os-lab10/xv6-riscv
Linux 6.14.0-37-generic
riscv64-linux-gnu-gcc (Ubuntu 13.3.0-6ubuntu2~24.04.1) 13.3.0
QEMU emulator version 8.2.2 (Debian 1:8.2.2+ds-0ubuntu1.18)
2441 README
\end{lstlisting}

The baseline built successfully with \texttt{make -j2}. A QEMU boot confirmed that
the original shell worked before any Lab 10 edits. QEMU was exited with its
Ctrl-A, X escape sequence after each checkpoint.

\begin{lstlisting}[caption={Initial boot and shell check; compiler output omitted}]
xv6 kernel is booting

hart 2 starting
hart 1 starting
init: starting sh
$ echo baseline boot OK
baseline boot OK
$ QEMU: Terminated
\end{lstlisting}

This version reserves syscall 22 for \texttt{sync} and uses \texttt{printk} for
kernel output, so the new syscall uses number 23 as required by the assignment.
The README is 2441 bytes, matching the assignment's expected read sequence.

\section{Part 1: The Number and User Side}
\subsection{Number, Declaration, Stub, and Disk Image}
The first stage adds \texttt{SYS\_trace} in \texttt{kernel/syscall.h}, declares
\texttt{int trace(int);} in \texttt{user/user.h}, and adds
\texttt{entry("trace");} to \texttt{user/usys.pl}. The Makefile's
\texttt{UPROGS} list receives a tab-indented \texttt{\$U/\_trace} entry followed
by a continuation backslash, placing the executable in the xv6 filesystem image.

The generated \texttt{user/usys.S} stub loads syscall number 23 into register
\texttt{a7}, executes \texttt{ecall}, and returns to the caller. The integer
argument is passed in \texttt{a0}. The following is an excerpt from the actual
stub-inspection output after the completed build:

\begin{lstlisting}[caption={Generated trace stub in user/usys.S}]
trace:
 li a7, SYS_trace
 ecall
 ret
\end{lstlisting}

\subsection{The User Command}
The file \texttt{user/trace.c} contains the program supplied by the assignment.
It validates the argument count, sets the current process's mask, and executes
the requested command. A successful \texttt{exec} replaces the program image
while retaining the process and its tracing mask.

\begin{lstlisting}[caption={Complete user/trace.c compiled in the VM}]
#include "kernel/types.h"
#include "kernel/stat.h"
#include "user/user.h"

int
main(int argc, char *argv[])
{
  if (argc < 3) {
    fprintf(2, "usage: trace mask command [args...]\n");
    exit(1);
  }
  if (trace(atoi(argv[1])) < 0) {
    fprintf(2, "trace: failed\n");
    exit(1);
  }
  exec(argv[2], argv + 2);
  fprintf(2, "trace: exec %s failed\n", argv[2]);
  exit(1);
}
\end{lstlisting}

\subsection{Check 1: Reaching an Unimplemented Kernel Call}
At this checkpoint only the user side was implemented. The tree was rebuilt and
booted before continuing to Part 2.

\begin{lstlisting}[caption={Check 1: actual expected failure before the handler is installed}]
$ trace 32 grep hello README
3 trace: unknown sys call 23
trace: failed
\end{lstlisting}

The unknown-call message proves that the program reached the kernel with number
23. The wrapper then saw a negative return value and printed \texttt{trace: failed}.
This is the intended intermediate result, not the result of the final implementation.

\section{Part 2: Kernel Implementation}
\subsection{Per-Process State and Handler}
An integer \texttt{trace\_mask} was added after the process name in
\texttt{struct proc}. The new \texttt{sys\_trace} handler obtains argument zero
with \texttt{argint}, saves the value on \texttt{myproc()}, and returns zero.

\begin{lstlisting}[caption={Handler added to kernel/sysproc.c}]
uint64
sys_trace(void)
{
  int mask;
  argint(0, &mask);
  myproc()->trace_mask = mask;
  return 0;
}
\end{lstlisting}

\subsection{Dispatch Table, Names, and Output}
The kernel declares \texttt{sys\_trace} and maps \texttt{SYS\_trace} to it in the
dispatch table. A table maps all syscall numbers from 1 through 23 to readable
names. After the selected handler returns, the dispatcher checks the corresponding
mask bit and prints the result already stored in \texttt{a0}:

\begin{lstlisting}[caption={Trace logic placed after syscall execution}]
p->trapframe->a0 = syscalls[num]();
if ((p->trace_mask >> num) & 1)
  printk("%d: syscall %s -> %ld\n",
         p->pid, syscallnames[num], p->trapframe->a0);
\end{lstlisting}

Placing this code after the handler is necessary to print the actual result.
It also means that a syscall which never returns, such as \texttt{exit}, does
not reach the post-call trace print. The kernel uses \texttt{printk} directly;
printing a trace line therefore does not recursively invoke a user-level
\texttt{write} system call.

\subsection{Initialization and Inheritance}
Newly allocated process state explicitly starts with a zero mask. When a parent
forks, the child's mask is copied immediately after its name. In this tree the
internal process-creation function is called \texttt{kfork}; it implements the
user-visible \texttt{fork} syscall.

\begin{lstlisting}[caption={Initialization and inheritance of the trace mask}]
// allocproc(), after setting the process state:
p->trace_mask = 0;

// kfork(), after copying the process name:
safestrcpy(np->name, p->name, sizeof(p->name));
np->trace_mask = p->trace_mask;
\end{lstlisting}

The explicit zero initialization supplements the assignment's edits with the
normal initialization of the added field. Inheritance allows descendants to
continue tracing. An \texttt{exec} does not overwrite this field, so the command
launched by the wrapper remains traced.

\subsection{Check 2: Exact Read Results}
The completed kernel built successfully and booted in QEMU. The required command
produced the following actual output:

\begin{lstlisting}[caption={Check 2: read-only tracing with the completed kernel}]
$ trace 32 grep hello README
3: syscall read -> 1023
3: syscall read -> 965
3: syscall read -> 453
3: syscall read -> 0
\end{lstlisting}

The successful reads total \(1023+965+453=2441\) bytes, matching README's size.
The final zero return value indicates end of file. Only reads appear because mask
32 is \(2^5\), and \texttt{SYS\_read} is 5. The staged verification checked this
exact sequence before proceeding to Part 3.

\section{Part 3: Mask Experiments}
\subsection{All Supported Calls}

\begin{lstlisting}[caption={All supported mask bits enabled}]
$ trace 2147483647 grep hello README
3: syscall trace -> 0
3: syscall exec -> 3
3: syscall open -> 3
3: syscall read -> 1023
3: syscall read -> 965
3: syscall read -> 453
3: syscall read -> 0
3: syscall close -> 0
\end{lstlisting}

The positive mask \(2147483647=2^{31}-1\) enables bits 0 through 30 and therefore
all syscall numbers currently used by this tree. The trace includes the
\texttt{trace} call itself because it installs the mask before the dispatcher
checks it. \texttt{exec} returns 3 in this xv6 implementation: the new program's
argument count. The open returns descriptor 3 and the close returns zero.

\subsection{Fork Only and Inherited Masks}
The required \texttt{forkforkfork} test was executed in a fresh QEMU session.
Below is the complete command output for the test, without the emulator startup
lines. Process identifiers and scheduling order can vary between executions.

\begin{lstlisting}[caption={Fork-only tracing and the required test result}]
$ trace 2 usertests forkforkfork
usertests starting
test forkforkfork: 3: syscall fork -> 4
4: syscall fork -> 5
5: syscall fork -> 6
5: syscall fork -> 7
6: syscall fork -> 8
6: syscall fork -> 9
7: syscall fork -> 10
5: syscall fork -> 11
6: syscall fork -> 12
6: syscall fork -> 13
5: syscall fork -> 14
5: syscall fork -> 15
6: syscall fork -> 16
6: syscall fork -> 17
6: syscall fork -> 18
5: syscall fork -> 19
6: syscall fork -> 20
5: syscall fork -> 21
6: syscall fork -> 22
6: syscall fork -> 23
6: syscall fork -> 24
5: syscall fork -> 25
6: syscall fork -> 26
5: syscall fork -> 27
5: syscall fork -> 28
6: syscall fork -> 29
6: syscall fork -> 30
5: syscall fork -> 31
6: syscall fork -> 32
5: syscall fork -> 33
5: syscall fork -> 34
5: syscall fork -> 35
5: syscall fork -> 36
6: syscall fork -> 37
5: syscall fork -> 38
6: syscall fork -> 39
6: syscall fork -> 40
6: syscall fork -> 41
5: syscall fork -> 42
6: syscall fork -> 43
6: syscall fork -> 44
6: syscall fork -> 45
5: syscall fork -> 46
6: syscall fork -> 47
6: syscall fork -> 48
6: syscall fork -> 49
47: syscall fork -> 50
5: syscall fork -> 51
5: syscall fork -> 52
6: syscall fork -> 53
5: syscall fork -> 54
48: syscall fork -> 55
5: syscall fork -> 56
6: syscall fork -> 57
5: syscall fork -> 58
6: syscall fork -> 59
5: syscall fork -> 60
6: syscall fork -> 61
21: syscall fork -> 62
5: syscall fork -> 63
6: syscall fork -> 64
5: syscall fork -> -1
6: syscall fork -> -1
OK
ALL TESTS PASSED
\end{lstlisting}

Multiple processes print fork results because children inherit the mask. The
negative return value \texttt{-1} represents a failed process allocation when
resources are exhausted. This is an expected part of the stress test. The test
completed with \texttt{OK} and \texttt{ALL TESTS PASSED}.

\begin{lstlisting}[caption={The ordinary shell remains usable after the fork test}]
$ echo after-fork-test
after-fork-test
\end{lstlisting}

\subsection{Write Only}

\begin{lstlisting}[caption={Write-only tracing, preserving the observed output interleaving}]
$ trace 65536 echo hi
hi4: syscall write -> 2

4: syscall write -> 1
\end{lstlisting}

Mask \(65536=2^{16}\) selects \texttt{SYS\_write}. The xv6 \texttt{echo} program
writes the two bytes of \texttt{hi} first, then writes a separate one-byte newline.
The kernel prints the first trace line immediately after the first write returns,
before \texttt{echo} issues its newline write. This is why the trace text follows
\texttt{hi} on the same line.

\subsection{Open and Close Only}
To select both calls, bit 15 and bit 21 must be set:
\[
  (1\ll15)\;|\;(1\ll21)=32768+2097152=2129920.
\]

\begin{lstlisting}[caption={Experimental verification of the open-and-close mask}]
$ trace 2129920 grep hello README
5: syscall open -> 3
5: syscall close -> 0
\end{lstlisting}

\subsection{Additional Wrapper Checks}
A zero mask suppressed trace lines while allowing the command to run. Missing
arguments displayed the usage message, and an unavailable command produced the
wrapper's execution-error message.

\begin{lstlisting}[caption={Additional command-line check}]
$ trace 0 echo untraced
untraced
\end{lstlisting}

\begin{lstlisting}[caption={Additional command-line check}]
$ trace
usage: trace mask command [args...]
\end{lstlisting}

\begin{lstlisting}[caption={Additional command-line check}]
$ trace 32 missingcommand
trace: exec missingcommand failed
\end{lstlisting}

\section{Observe Questions and Answers}
The answers below are also saved in both \texttt{report.md} and the assignment's
requested \texttt{notes.txt} under \texttt{/home/user/os-lab10}.

\subsection{Why Only Read Calls with Mask 32?}
\textbf{Question:} In \texttt{trace 32 grep hello README}, why do you see only
\texttt{read}, and not \texttt{open} or \texttt{close}?

\textbf{Answer:} The binary representation of 32 has only bit 5 set. That bit
selects syscall 5, \texttt{read}. The bits for \texttt{open} (15) and
\texttt{close} (21) are clear. Those operations still execute; their trace lines
are suppressed.

\subsection{Why Do Descendants Print Trace Lines?}
\textbf{Question:} In \texttt{forkforkfork}, different process identifiers print
lines. Which line from Part 2 makes that possible?

\textbf{Answer:} The assignment's inheritance line is
\texttt{np->trace\_mask = p->trace\_mask;}. It copies the parent's mask into the
child during process creation, so the child's own fork calls are also traced.

\subsection{Why Does a Trace Line Interrupt the Echo Output?}
\textbf{Question:} Why is the trace line between \texttt{hi} and its newline?
How many times does \texttt{echo} call \texttt{write}?

\textbf{Answer:} For \texttt{echo hi}, there are two writes: one of two bytes
for \texttt{hi} and one of one byte for the newline. The kernel prints each
trace line after its corresponding write returns. The first trace therefore
appears before the program has written its newline.

\subsection{Which Mask Selects Open and Close?}
\textbf{Question:} Which mask traces \texttt{open} (15) and \texttt{close} (21),
and nothing else?

\textbf{Answer:} \texttt{2129920}, the bitwise OR of \texttt{1 << 15} and
\texttt{1 << 21}. The experiment above confirms that it prints only an open
and a close for the tested command.

\section{Results and Committed Deliverable}
The initial tree, user-side checkpoint, and final kernel each built and booted.
Both required checkpoints matched the expected behavior. The fork-inheritance
test passed, and all three required mask experiments were completed. Additional
checks verified the combined open/close mask, zero-mask behavior, and wrapper
errors. \texttt{git diff --check} reported no whitespace problems.

The completed implementation is committed in the VM working tree on branch
\texttt{codex/lab10-trace}. The following output records the final commit and
changed files:

\begin{lstlisting}[caption={Recorded Git history, working-tree status, and implementation summary}]
$ bash -c git log -2 --oneline; git status --short; git show --stat --oneline HEAD
4f95820 Add per-process trace system call and mask inheritance
3952f44 Preserve existing coursework before Lab 10
?? args
4f95820 Add per-process trace system call and mask inheritance
 Makefile         |  1 +
 kernel/proc.c    |  2 ++
 kernel/proc.h    |  1 +
 kernel/syscall.c | 31 +++++++++++++++++++++++++++++++
 kernel/syscall.h |  1 +
 kernel/sysproc.c |  9 +++++++++
 user/trace.c     | 19 +++++++++++++++++++
 user/user.h      |  1 +
 user/usys.pl     |  1 +
 9 files changed, 66 insertions(+)
\end{lstlisting}

The untracked \texttt{args} file is a pre-existing compiled program copied from
the earlier coursework tree; it is not part of the Lab 10 implementation commit.
All nine implementation files, including the new \texttt{user/trace.c}, are
committed. A portable Git bundle containing the complete history and this branch
was saved as \texttt{xv6-lab10.bundle} in the VM lab directory and copied to the
host's \texttt{Reports/source\_code\_and\_materials\_of\_all\_labs/os-lab10}
folder. The bundle was verified successfully on the host. It can be restored with:

\begin{lstlisting}[caption={Restoring the committed xv6 source from the portable Git bundle}]
git clone -b codex/lab10-trace xv6-lab10.bundle xv6-riscv
\end{lstlisting}

This restoration command is provided for future use; it is not a recorded test
session. The original VM working tree remains ready for \texttt{make qemu}.

\section{Conclusion}
The new \texttt{trace} system call connects a user-level request to a kernel
handler through a generated RISC-V stub and the syscall dispatch table. Per-process
mask state selects the events to print, persists across \texttt{exec}, and is
inherited by children. The successful read sequence, fork stress test, and write
interleaving demonstrate that the implementation behaves as required. The work
also illustrates why tracing must occur after a handler returns and why syscall
return values and process scheduling affect the observed output.

\appendix
\section{Recorded Implementation Changes}
The following is the actual diff against the preserved pre-Lab-10 coursework
baseline for the eight existing files. The complete new \texttt{user/trace.c}
file is reproduced in Part 1. Together they document all nine implementation files.

\begin{lstlisting}[caption={Actual changes to existing xv6 files}]
$ git diff -- Makefile kernel/proc.h kernel/proc.c kernel/sysproc.c kernel/syscall.h kernel/syscall.c user/user.h user/usys.pl
diff --git a/Makefile b/Makefile
index 3accbe6..2813348 100644
--- a/Makefile
+++ b/Makefile
@@ -133,6 +133,7 @@ UPROGS=\
 	$U/_cat\
 	$U/_find\
 	$U/_zombie\
+	$U/_trace\
 	$U/_run\
 	$U/_pingpong\
 	$U/_echo\
diff --git a/kernel/proc.c b/kernel/proc.c
index 31f23a3..2dd5d14 100644
--- a/kernel/proc.c
+++ b/kernel/proc.c
@@ -124,6 +124,7 @@ allocproc(void)
 found:
   p->pid = allocpid();
   p->state = USED;
+  p->trace_mask = 0;
 
   // Allocate a trapframe page.
   if ((p->trapframe = (struct trapframe *)kalloc()) == 0) {
@@ -288,6 +289,7 @@ kfork(void)
   np->cwd = idup(p->cwd);
 
   safestrcpy(np->name, p->name, sizeof(p->name));
+  np->trace_mask = p->trace_mask;
 
   pid = np->pid;
 
diff --git a/kernel/proc.h b/kernel/proc.h
index 5603697..121d8ad 100644
--- a/kernel/proc.h
+++ b/kernel/proc.h
@@ -101,4 +101,5 @@ struct proc {
   struct file *ofile[NOFILE];  // Open files
   struct inode *cwd;           // Current directory
   char name[16];               // Process name (debugging)
+  int trace_mask;              // Bit N selects syscall N
 };
diff --git a/kernel/syscall.c b/kernel/syscall.c
index 1800f20..8df9614 100644
--- a/kernel/syscall.c
+++ b/kernel/syscall.c
@@ -103,6 +103,7 @@ extern uint64 sys_link(void);
 extern uint64 sys_mkdir(void);
 extern uint64 sys_close(void);
 extern uint64 sys_sync(void);
+extern uint64 sys_trace(void);
 
 // An array mapping syscall numbers from syscall.h
 // to the function that handles the system call.
@@ -130,9 +131,36 @@ static uint64 (*syscalls[])(void) = {
   [SYS_mkdir]   = sys_mkdir,
   [SYS_close]   = sys_close,
   [SYS_sync]    = sys_sync,
+  [SYS_trace]   = sys_trace,
   // clang-format on
 };
 
+static char *syscallnames[] = {
+  [SYS_fork] = "fork",
+  [SYS_exit] = "exit",
+  [SYS_wait] = "wait",
+  [SYS_pipe] = "pipe",
+  [SYS_read] = "read",
+  [SYS_kill] = "kill",
+  [SYS_exec] = "exec",
+  [SYS_fstat] = "fstat",
+  [SYS_chdir] = "chdir",
+  [SYS_dup] = "dup",
+  [SYS_getpid] = "getpid",
+  [SYS_sbrk] = "sbrk",
+  [SYS_pause] = "pause",
+  [SYS_uptime] = "uptime",
+  [SYS_open] = "open",
+  [SYS_write] = "write",
+  [SYS_mknod] = "mknod",
+  [SYS_unlink] = "unlink",
+  [SYS_link] = "link",
+  [SYS_mkdir] = "mkdir",
+  [SYS_close] = "close",
+  [SYS_sync] = "sync",
+  [SYS_trace] = "trace",
+};
+
 void
 syscall(void)
 {
@@ -144,6 +172,9 @@ syscall(void)
     // Use num to lookup the system call function for num, call it,
     // and store its return value in p->trapframe->a0
     p->trapframe->a0 = syscalls[num]();
+    if ((p->trace_mask >> num) & 1)
+      printk("%d: syscall %s -> %ld\n",
+             p->pid, syscallnames[num], p->trapframe->a0);
   } else {
     printk("%d %s: unknown sys call %d\n", p->pid, p->name, num);
     p->trapframe->a0 = -1;
diff --git a/kernel/syscall.h b/kernel/syscall.h
index 5a4a39f..f66fb28 100644
--- a/kernel/syscall.h
+++ b/kernel/syscall.h
@@ -21,3 +21,4 @@
 #define SYS_mkdir  20
 #define SYS_close  21
 #define SYS_sync   22
+#define SYS_trace  23
diff --git a/kernel/sysproc.c b/kernel/sysproc.c
index cd25557..f3b5d40 100644
--- a/kernel/sysproc.c
+++ b/kernel/sysproc.c
@@ -110,3 +110,12 @@ sys_uptime(void)
   release(&tickslock);
   return xticks;
 }
+
+uint64
+sys_trace(void)
+{
+  int mask;
+  argint(0, &mask);
+  myproc()->trace_mask = mask;
+  return 0;
+}
diff --git a/user/user.h b/user/user.h
index e571a5e..66ad64a 100644
--- a/user/user.h
+++ b/user/user.h
@@ -26,6 +26,7 @@ char *sys_sbrk(int, int);
 int pause(int);
 int uptime(void);
 int sync(void);
+int trace(int);
 
 // ulib.c
 int stat(const char *, struct stat *);
diff --git a/user/usys.pl b/user/usys.pl
index 4d054e8..f9185b5 100755
--- a/user/usys.pl
+++ b/user/usys.pl
@@ -43,3 +43,4 @@ entry("sbrk");
 entry("pause");
 entry("uptime");
 entry("sync");
+entry("trace");
\end{lstlisting}

\end{document}
